\chapter{相对论}\label{chap:relativity}
\begingroup
\renewcommand{\thefootnote}{}
\footnotetext{\textcopyright\ Springer International Publishing Switzerland 2014.
  L. Godinho and J. Natário, \emph{An Introduction to Riemannian Geometry},
  Universitext. DOI \href{https://doi.org/10.1007/978-3-319-08666-8_6}{10.1007/978-3-319-08666-8\_6}.}
\endgroup

本章研究 Riemann 几何最重要的应用之一: \textbf{广义相对论}. 这一理论最终取代了 Galileo 和 Newton 的经典力学, 它源于一个看似矛盾的实验事实: 光速对每个观察者都是相同的, 与观察者的运动状态无关. 1905 年, 经历了一段极为困惑的时期后, Einstein 提出了一个既简单又彻底的解释: 时间间隔和长度测量值并不对所有观察者相同, 而是取决于各自的运动状态. 1908 年, Minkowski 在四维\textbf{时空} $\mathbb{R}^4$ 中引入伪内积, 给出了 Einstein 理论的一种几何表述. Einstein 起初抵制对其理论的这种\textquotedblleft 过度数学化\textquotedblright, 但很快认识到, 时空弯曲实际上是理解引力的关键. 1915 年, 在长期钻研 Riemann 几何的数学之后, 他终于给出了广义相对论的完整表述. 他的理论预言于 1919 年首先得到 Eddington 率领的英国日食观测考察队的证实, 此后也在每一次进行过的实验检验中得到验证.

为了从经典力学自然过渡到相对论, 第~\ref{sec:galileo-spacetime}~节讨论 \textbf{Galileo 时空}, 即作为 Newton 力学基础的几何结构; 这一结构有赖于任意高速运动的存在. 然而, 若假定存在最大速度, 则必须以 \textbf{Minkowski 时空}取代它, 其几何由\textbf{狭义相对论}研究 (第~\RGref{sec:special-relativity}~节).

第~\RGref{sec:cartan-connection}~节通过引入对称的 \textbf{Cartan 联络}, 说明如何将 Newton 引力纳入 Galileo 时空. 试图将这一过程推广, 会得到满足 \textbf{Einstein 场方程}的一般 \textbf{Lorentz 流形}, 而 Minkowski 时空是其中最简单的例子 (第~\RGref{sec:general-relativity}~节).

后续各节分析其他简单的解: 用来模拟球对称物体或\textbf{黑洞}外部引力场的 \textbf{Schwarzschild 解} (第~\RGref{sec:schwarzschild-solution}~节), 以及\textbf{宇宙学}中的 \textbf{Friedmann–Lemaître–Robertson–Walker 模型}, 它们描述宇宙的整体行为 (第~\RGref{sec:cosmology}~节).

最后, 第~\RGref{sec:causality-in-lorentzian-manifolds}~节讨论 Lorentz 流形的\textbf{因果结构}, 为证明著名的 \textbf{Hawking} (第~\RGref{sec:hawking-singularity-theorem}~节) 和 \textbf{Penrose} (第~\RGref{sec:penrose-singularity-theorem}~节) \textbf{奇性定理}作准备.
